\documentclass[a4paper,12pt]{article}
\setlength{\headheight}{68.803pt}
\addtolength{\topmargin}{-56.803pt}

\usepackage[utf8]{inputenc}
\usepackage[T1]{fontenc}
\usepackage{lmodern}
\usepackage[nopatch=footnote]{microtype}
\usepackage{fvextra}
\DefineVerbatimEnvironment{Verbatim}{Verbatim}{
  breaklines=true,
  breakanywhere=true,
  frame=single,
  fontsize=\small
}

\usepackage[
  left=1in,
  right=1in,
  top=3cm,
  bottom=1in
]{geometry}

\usepackage[
    colorlinks=true,
    linkcolor=black,
    urlcolor=blue,
    citecolor=blue,
    pdfborder={0 0 0},
    pdfauthor={Massimo Mantineo},
    pdftitle={HO15: Flooding con FIFO},
    pdfsubject={Laboratorio di Reti e Sistemi Distribuiti: HANDS-ON 15},
]{hyperref}

\usepackage{graphicx}
\graphicspath{{../../assets/}}
\usepackage{tikz}
\usetikzlibrary{positioning, arrows.meta, shapes.symbols}
\usepackage{listings}
\usepackage{xcolor}
\usepackage{amsmath, amssymb}
\usepackage{fancyhdr}
\usepackage{setspace}
\usepackage{pgfplots}
\pgfplotsset{compat=1.17}
\usepackage{booktabs}
\usepackage[skins, listings]{tcolorbox}
\usepackage{float}

\newtcblisting{terminal}[1][]{
    listing only,
    colback=black!85!white,
    colframe=black!70!white,
    colupper=green!80!white,
    coltitle=white,
    title=#1,
    fonttitle=\bfseries\sffamily,
    listing options={
        language=bash,
        basicstyle=\ttfamily\small\color{green!80!white},
        backgroundcolor={},
        frame=none,
        breaklines=true,
        showstringspaces=false,
        numbers=none
    },
    enhanced,
    drop shadow,
    arc=3mm
}

\newtcblisting{ccode}[1][]{
    listing only,
    colback=blue!3!white,
    colframe=blue!70!black,
    title=#1,
    fonttitle=\bfseries\sffamily,
    listing options={
        style=cstyle,
        backgroundcolor={},
        frame=none
    },
    enhanced,
    drop shadow,
    arc=3mm
}

\setlength{\footskip}{50pt}

\lstdefinestyle{cstyle}{
    language=C,
    basicstyle=\ttfamily\small,
    numbers=left,
    numberstyle=\tiny,
    stepnumber=1,
    numbersep=5pt,
    backgroundcolor=\color{white},
    showspaces=false,
    showstringspaces=false,
    showtabs=false,
    frame=single,
    tabsize=4,
    captionpos=b,
    breaklines=true,
    breakatwhitespace=true,
    keywordstyle=\color{blue}\bfseries,
    commentstyle=\color{green!50!black},
    stringstyle=\color{red},
}
\lstset{style=cstyle}

\pagestyle{fancy}
\fancyhf{}

\fancyhead[L]{%
  \parbox[b]{0.45\textwidth}{%
    \small
    Student Name: Massimo Mantineo\\
    Student ID: 541924
    \vspace{20pt}
  }%
}
\fancyhead[C]{%
  \parbox[b]{0.2\textwidth}{%
    \centering
    \normalsize \bfseries
    LabReSiD25\\
    Hands-On 15
    \vspace{20pt}
  }%
}
\fancyhead[R]{%
  \parbox[b]{0.35\textwidth}{%
    \flushright
    \includegraphics[width=3.5cm]{\detokenize{logo_unime_orizontale.png}}
    \vspace{15pt}
  }%
}

\renewcommand{\contentsname}{Indice}

\title{\vspace{-2cm}Relazione (HO15)\\
\large Algoritmo Generico di Flooding via FIFO / IPC Processes}
\author{Massimo Mantineo \\ Matricola: 541924}
\date{MESSINA, A.A. 2024/2025}

\begin{document}

\begin{titlepage}
    \thispagestyle{empty}
    \centering
    \vspace*{1cm}
    \includegraphics[width=6cm]{\detokenize{logo_unime_orizontale.png}}\par\vspace{1cm}
    \vspace{2cm}
    {\Large \textbf{Università degli Studi di Messina}\par}
    {\large Dipartimento MIFT - Corso di Laurea Triennale in Informatica (L-31)\par}
    \vspace*{1cm}
    {\large Laboratorio di Reti e Sistemi Distribuiti\par}
    \vspace{1cm}
    {\Large \textbf{HANDS-ON 15:}\\
    Implementazione del Flooding Algorithm su Processi POSIX\par}
    \vspace{2cm}
    {\large Massimo Mantineo \par}
    {\large Matricola: 541924 \par}
    \vfill
    {\large Messina, A.A. 2024/2025 \par}
\end{titlepage}

\clearpage
\pagestyle{fancy}
\tableofcontents
\newpage

\section{Introduzione al Problema}
L'Hands-On 15 richiede la dismissione dei costrutti standard d'Ateneo tramite \texttt{Pthreads} o i classici Web Sockets TCP, allo scopo di studiare la comunicazione Inter-Processo locale di Linux (IPC) tramite le \textbf{Named Pipes} (conosciute come FIFO). L'applicazione designata è la programmazione distribuita di un algoritmo "Flooding", in cui viaggiano messaggi da un Iniziatore (Agente) all'insegna di vicini collegati direzionalmente su un grafo, finchè tutta la rete non è informata. A un nodo (\textit{Processo isolato generato da una Fork()}) e alle comunicazioni sui suoi archi fisici (\textit{Le FIFO in /tmp/}) dovranno asserragliarsi le funzioni \texttt{msg()} ed \texttt{stf()}.

\section{Stato dell'Arte}
A livello concettuale, le IPC Pipelines vengono comunemente introdotte da Ritchie. 
\begin{itemize}
    \item \textbf{Anonymous vs Named Pipes}: A differenza di `pipe(fd)`, i `mkfifo()` esistono nel Filesystem Linux e possono essere ispezionati ed aperti dai binari isolatamente solo per nome.
    \item \textbf{Modello a Nodi / Agenti}: Isolare la computazione usando `fork()` è imperativo in un lab sui grafi, poiché la memoria separata previene trucchi di lettura "illegale" (come array globali condivisi ai thread). L'unico modo in cui il Nodo 2 può sapere cosa ha pensato il Nodo 0 è aspettare sulla propria Pipeline.
    \item \textbf{Costrutti Formali}: L'algoritmo di Santoro identifica il Message Generator ed il valutatore interno (State Transition Function, \texttt{stf}).
\end{itemize}

\section{Metodologia usata}

\subsection{Topologia ed Avvio}
Trattandosi di uno scenario aperto, abbiamo optato per formalizzare un tipico **Grafo a Diamante di 4 Nodi**: L'Agent 0 (Iniziatore) è connesso asimmetricamente ad {1, 2}. Gli Agent 1 e 2 si rivolgono entrambi al nodo di base {3}.

Dal parent \texttt{main()} del C, il master lancia ciclicamente la chiamata di Sistema \texttt{mkfifo("/tmp/fifo\_src\_dst", 0666)} generando i tubi direzionali:
\begin{ccode}[Creazione Dinamica]
for(int i = 0; i < MAX_AGENTS; i++) {
    for(int j = 0; j < num_neighbors[i]; j++) {
        char fname[64];
        sprintf(fname, "/tmp/fifo_%d_%d", i, neighbors[i][j]);
        mkfifo(fname, 0666);
    }
}
\end{ccode}
Dopo \texttt{fork()} parallela, ogni figlio agisce come Agente e non può più parlare con il fratello se non usando la pipeline.

\subsection{stf() e Multiplexing}
Le FIFO in ambiente Linux subiscono la natura \textit{Bloccante} della Syscall \texttt{open(... , O\_RDONLY)}. Quando si chiamano multipli Open, se il vicino non avvia la controparte Write il Processo freezza. È stata utilizzata la bandiera \texttt{O\_RDWR} (che Linux autorizza per risolvere i deadlock di bootstrap) per poi gettare i rispettivi FD dentro un \texttt{fd\_set} processato da \texttt{select()}.

\begin{ccode}[State Transition Function]
void stf(int id, int* state, int sender, const char* info) {
    if (*state == 0) { // PASSIVE TO DONE
        *state = 1; 
        for(int i = 0; i < num_neighbors[id]; i++) {
            int neigh = neighbors[id][i];
            if (neigh != sender) msg(id, neigh, info);
        }
    }
}
\end{ccode}

\section{Risultati Sperimentali}
Avviando la Routine, la stampa al terminale esalta coerentemente la logica asincrona dei 4 nodi isolati e del Flooding, in cui i Nodi 2 e 3 alla fine scartano i messaggi ritardatari pervenuto dell'altro braccio.

\begin{terminal}
massimo@PORTATILE-Linux:~/Scrivania/LabReSid24-25/hands-on/ho15$ ./flooding
--- Avvio Algoritmo di Flooding Centralizzato ---
Topologia: 4 Nodi a Diamante (0 Initiator)

=> [AGENTE 0] INIZIATORE: Genero il messaggio di Flooding iniziale.
    -> [AGENTE 0] Propaga messaggio al Nodo 1
    -> [AGENTE 0] Propaga messaggio al Nodo 2
[AGENTE 1] RICEZIONE da Nodo 0: 'ALLERTA_METEO'. TRASIZIONE a DONE e Propagazione.
    -> [AGENTE 1] Propaga messaggio al Nodo 3
[AGENTE 2] RICEZIONE da Nodo 0: 'ALLERTA_METEO'. TRASIZIONE a DONE e Propagazione.
    -> [AGENTE 2] Propaga messaggio al Nodo 3
[AGENTE 3] RICEZIONE da Nodo 1: 'ALLERTA_METEO'. TRASIZIONE a DONE e Propagazione.
    -> [AGENTE 3] Propaga messaggio al Nodo 2
[AGENTE 2] SCARTO msg da Nodo 3 (Gia' elaborato).
[AGENTE 3] SCARTO msg da Nodo 2 (Gia' elaborato).
[AGENTE 2] Spegnimento.
[AGENTE 1] Spegnimento.
[AGENTE 3] Spegnimento.
[AGENTE 0] Spegnimento.
\end{terminal}

A fine esecuzione, il master rimuove ordinatamente dalle cartelle UNIX \texttt{/tmp/} la sequela di FIFO in disuso per mantenere il Filesystem esente da clutters.

\section{Conclusione e Sviluppi Futuri}
L'algoritmo riproduce con precisione estrema un albero decisionale DFS-like senza attese esterne se non quelle bloccanti native delle pipe POSIX.

\noindent \textbf{Sviluppi futuri ideali:}
\begin{itemize}
    \item \textbf{Polimorfismo}: Leggere tramite argomento iniziale il layout di un Grafo qualsiasi esportato dal Software NodeRED o passati via struct JSON per dinamicizzare l'orchestratore, anziché codificare la topologia esplicitamente nel C.
    \item \textbf{Messaggistica Bidirezionale pura}: Usare protocolli a coda per stabilire il Teardown anziché usare un \textit{Timeout empirico} (Il nodo termina dopo N cicli morti dalla \texttt{select}). Si eviterebbero cicli clock rubati a vuoto dallo scheduler OS.
\end{itemize}

\end{document}
